\section{A uniform frame from binned Dirichlet zero counts}

The new point here is uniformity in every real coefficient vector and in every
microscopic partition.  The binary Hellinger estimate of order $T\log T$ is
already present in the incoming report
\texttt{cycle-complements-and-dirichlet-zero-matching/03-hard-distance-matching}.
It is not claimed as new here~\cite{BinaryIncoming}.  The frame theorem below extends that report
to continuously varying mixtures and yields a minimax estimation theorem.

Let $\chi_1,\ldots,\chi_d$ be distinct fixed primitive Dirichlet characters,
including the conductor-one character if desired.  Count positive ordinates
of nontrivial zeros with multiplicity, without requiring the zeros to lie on
a vertical line.  Put $\mathcal L_T=T\log T$.  Let $\mathcal P_T$ be a
partition of $(T,2T]$ into measurable bins, each of diameter at most
$K/\log T$, where $K<\infty$ is fixed.  Empty bins can be removed, leaving
finitely many bins that contain a zero of the family.  Define
\[
 v_B=(N_1(B),\ldots,N_d(B))^{\mathsf T},\qquad
 s_B=\boldsymbol 1^{\mathsf T}v_B,
 \qquad
 G_T=\sum_{B:s_B>0}\frac{v_Bv_B^{\mathsf T}}{s_B}.
\]

\begin{theorem}[Uniform microscopic frame]\label{mix:thm:frame}
There are $c,C,T_0>0$, depending only on the fixed character family and $K$,
such that for $T\ge T_0$, uniformly over all such partitions,
\[
 c\mathcal L_T I_d\preceq G_T\preceq C\mathcal L_T I_d.
\tag{MF1}
\]
This is an unconditional theorem about zero ordinates.  No assertion about
their simplicity, separation, or noncoincidence is needed.
\end{theorem}

\subsection{The classical input}

Write $C_j(u,h)=N_j(u+h)-N_j(u)$ and let $u$ be uniform on $[T,2T]$.
The short-interval joint central limit theorem of Bombieri and Perelli
\cite[Theorem 2]{BP1998} implies that whenever
\[
 M=M(T)\longrightarrow\infty,\qquad
 M\le\frac{\log T}{\log\log T},\qquad h=M/\log T,
\]
the vector
\[
 Z_T=\left(\frac{C_j(u,h)-M/(2\pi)}{\sqrt{\log M}}\right)_{j=1}^d
\tag{MF2}
\]
converges in distribution to a centered nondegenerate Gaussian vector;
indeed its covariance is $\pi^{-2}I_d$.  The Riemann--von Mangoldt formula
absorbs the fixed conductor terms and the replacement of $\log u$ by
$\log T$ into $O(h+T^{-1})$.

For completeness, the hypotheses in that paper have established
unconditional verification for this family.  The Euler parameters have
modulus at most one, the continuations and functional equations are
classical, and character orthogonality plus the prime number theorem in
fixed residue classes gives their prime orthogonality hypothesis B(IV).
Their density hypothesis D is known for Dirichlet $L$-functions, as explicitly
stated on page 273 of the paper.  Thus (MF2) is not being imported from an
unchecked assertion in the research article.  For $d=1$, one may adjoin a
different primitive character and take a marginal.

We will also use the first-moment consequence
\[
 \mathbb E|C_j(u,h)-C_k(u,h)|=O\bigl(\sqrt{\log M}\bigr)
 \qquad(j,k\text{ fixed}).
\tag{MF3}
\]
This does not follow from weak convergence alone.  It follows from the
$L^1$ approximation by normalized Dirichlet polynomials on page 279 of
\cite{BP1998}, together with the bounded second moments in their Lemma 2.
The fixed conductor drift is $O(h)$ and is negligible.

\subsection{Proof of the frame theorem}

First we prove the uniform $\ell^1$ lower bound
\[
 \sum_B|b^{\mathsf T}v_B|\ge c_1\mathcal L_T\|b\|_2
 \qquad(b\in\mathbb R^d).
\tag{MF4}
\]
Suppose it fails.  Then along a sequence $T\to\infty$ there are allowed
partitions and unit vectors $b_T$ such that
\[
 D_T:=\sum_B|b_T^{\mathsf T}v_B|=o(\mathcal L_T).
\]
Pass to a subsequence with $b_T\to b$, $\|b\|_2=1$, and write
$\delta_T=D_T/\mathcal L_T$.  Choose
\[
 M=\min\left\{\frac{\log T}{\log\log T},
        (\delta_T+\mathcal L_T^{-1})^{-1/2}\right\}.
\]
Then $M\to\infty$ and $M\delta_T\to0$.

Choose one representative $x_B$ in each nonempty bin and move every zero
occurrence in that bin to $x_B$.  Each move has length at most
$K/\log T$.  For a real $x$, the indicator that $x$ belongs to $(u,u+h]$
is supported on the interval $[x-h,x)$ of starting points.  Moving $x$
to $x'$ changes this indicator on a set of measure
$2\min(h,|x-x'|)$.  Therefore, on $u\in[T,2T-h]$,
\begin{align*}
 &\int_T^{2T-h}\left|\sum_j b_{T,j}C_j(u,h)
      -\sum_B b_T^{\mathsf T}v_B\,
                    \boldsymbol 1_{(u,u+h]}(x_B)\right|\,du
 \\
 &\hspace{25mm}\le
 \frac{2K}{\log T}\sum_j|b_{T,j}|\bigl(N_j(2T)-N_j(T)\bigr)
 =O(KT).
\end{align*}
The integral of the absolute value of the second sum is at most $hD_T$.
Thus
\[
 \frac1T\int_T^{2T-h}\left|\sum_j b_{T,j}C_j(u,h)\right|du
 \le M\delta_T+O(K).
\tag{MF5}
\]
The discarded terminal interval of starting points has probability
$h/T=o(1)$.  Divide (MF5) by $\sqrt{\log M}$.  It says that
\[
 a_T+b_T^{\mathsf T}Z_T\longrightarrow0
 \quad\hbox{in probability},\qquad
 a_T=\frac{M\sum_jb_{T,j}}{2\pi\sqrt{\log M}}.
\]
But $b_T^{\mathsf T}Z_T$ converges to a nondegenerate normal law.  If $a_T$
has a bounded subsequence, take a further convergent subsequence to get a
nondegenerate translated normal limit, a contradiction.  If $|a_T|\to\infty$
along a subsequence, tightness of $b_T^{\mathsf T}Z_T$ is a contradiction.
This proves (MF4).

The total mass satisfies $\sum_Bs_B=O(\mathcal L_T)$.  Cauchy--Schwarz gives
\[
 \left(\sum_B|b^{\mathsf T}v_B|\right)^2
 \le\left(\sum_Bs_B\right)
      \left(\sum_B\frac{(b^{\mathsf T}v_B)^2}{s_B}\right),
\]
which proves the lower bound in (MF1).  For the upper bound use weighted
Cauchy--Schwarz separately in every bin:
\[
 \frac{(b^{\mathsf T}v_B)^2}{s_B}
 \le\sum_j b_j^2(v_B)_j.
\]
After summing, each diagonal total is $O(\mathcal L_T)$ by zero counting.
This completes the proof.

\section{Optimal recovery of continuous mixtures}

Suppose independently for each bin we observe
\[
 Y_B\sim\operatorname{Poisson}(t\,\theta^{\mathsf T}v_B),
 \qquad \theta\in[0,R]^d,
\tag{MF6}
\]
where $t>0$ is known and the count vectors $v_B$ are known.  The parameter
$\theta$ specifies an intensity superposition, not a random choice of one
component.  It need not have integer or rational entries.

\begin{theorem}[Minimax rate and an explicit estimator]\label{mix:thm:estimator}
With the setup of Theorem~\ref{mix:thm:frame}, define
\[
 \widehat\theta=\frac1tG_T^{-1}
                  \sum_{B:s_B>0}\frac{v_B}{s_B}Y_B.
\tag{MF7}
\]
Then $\mathbb E_\theta\widehat\theta=\theta$ and
\[
 \operatorname{Cov}_\theta(\widehat\theta)
 =\frac1tG_T^{-1}\left[
     \sum_B\frac{\theta^{\mathsf T}v_B}{s_B^2}v_Bv_B^{\mathsf T}
                       \right]G_T^{-1}
 \preceq\frac RtG_T^{-1}.
\tag{MF8}
\]
If $G_T\succeq c\mathcal L_T I_d$, then for every $u>0$,
\[
 \Pr_\theta\{\|\widehat\theta-\theta\|_\infty\ge u\}
 \le2d\exp\left\{-\frac{ct\mathcal L_Tu^2}{2R+2u/3}\right\}.
\tag{MF9}
\]
Projecting (MF7) onto $[0,R]^d$ gives an estimator whose worst-case squared
Euclidean risk is $O(\min\{1,(t\mathcal L_T)^{-1}\})$.  Conversely, for every
fixed $0<r<R$, even the full unbinned experiment restricted to this band satisfies
\[
 \inf_{\widetilde\theta}\sup_{\theta\in[r,R]^d}
 \mathbb E_\theta\|\widetilde\theta-\theta\|_2^2
 \ge c_{r,R}\min\{1,(t\mathcal L_T)^{-1}\}.
\tag{MF10}
\]
All constants are uniform over the allowed bin partitions.  They depend on
the fixed family, $K,d,r,R$; no conductor-uniform assertion is made.
\end{theorem}

\begin{proof}
The expectation in (MF7) is
$G_T^{-1}\sum_Bv_Bv_B^{\mathsf T}\theta/s_B=\theta$.
Independence and the Poisson variance formula give the equality in (MF8),
and $\theta^{\mathsf T}v_B\le Rs_B$ gives its inequality.

For a fixed coordinate $j$, put
$q_B=e_j^{\mathsf T}G_T^{-1}v_B/s_B$.  Since $v_B$ has nonnegative
coordinates, $\|v_B\|_2\le s_B$, so $|q_B|\le1/(c\mathcal L_T)$.
The centered sum $t(\widehat\theta_j-\theta_j)$ has log moment generating
function
$t\sum_B(\theta^{\mathsf T}v_B)(e^{zq_B}-1-zq_B)$.
The bound
$e^x-1-x\le x^2/[2(1-|x|/3)]$ for $|x|<3$, followed by exponential
Markov, gives a two-sided Bernstein bound with variance parameter at most
$tR/(c\mathcal L_T)$ and maximum jump size $1/(c\mathcal L_T)$.
Substitution and a union bound yield (MF9).  Taking the trace in (MF8)
gives the stated risk.  Coordinatewise clipping cannot increase squared
error relative to any parameter in the box, and bounds it by $dR^2$.

For the lower bound, take two interior points differing by
$\delta e_1$, with $\delta$ smaller than a fixed fraction of $R-r$.
For Poisson laws with means $a,b>0$,
\[
 D_{\rm KL}(\operatorname{Pois}(a)\,\|\,\operatorname{Pois}(b))
 =a\log(a/b)-a+b\le(a-b)^2/b.
\]
Thus, including in the unbinned band experiment,
\[
 D_{\rm KL}(P_\theta\|P_\eta)
 \le\frac{t\delta^2}{r}\sum_B\frac{(v_B)_1^2}{s_B}
 \le\frac{t\delta^2}{r}\sum_B(v_B)_1
 \le C t\mathcal L_T\delta^2.
\]
Choose $\delta$ a sufficiently small constant times
$\min\{1,(t\mathcal L_T)^{-1/2}\}$.  Pinsker's inequality makes the total
variation distance at most $1/2$.  The two-point testing bound for squared
loss then gives (MF10).
\end{proof}

\paragraph{Fisher information and Hellinger geometry.}
For $\theta\in[r,R]^d$ the exact Fisher information matrix is
\[
 \mathcal I_T(\theta)
 =t\sum_B\frac{v_Bv_B^{\mathsf T}}{\theta^{\mathsf T}v_B},
 \qquad
 \frac tR G_T\preceq\mathcal I_T(\theta)\preceq\frac trG_T.
\tag{MF11}
\]
For $\theta,\eta\in[r,R]^d$ the exact squared-root intensity discrepancy
satisfies
\[
 \frac{(\theta-\eta)^{\mathsf T}G_T(\theta-\eta)}{4R}
 \le \sum_B(\sqrt{\theta^{\mathsf T}v_B}
                 -\sqrt{\eta^{\mathsf T}v_B})^2
 \le\frac{(\theta-\eta)^{\mathsf T}G_T(\theta-\eta)}{4r}.
\tag{MF12}
\]
Both inequalities follow by rationalizing each square-root difference.
The Poisson affinity is the exponential of $-t/2$ times this sum.
The new content is the uniform nondegenerate matrix bound, rather than
the standard affinity identity.

\section{Exactly \texorpdfstring{$d$}{d} aggregated Poisson counters suffice}

The estimator (MF7) reduces a large histogram to a $d$-dimensional weighted
statistic, but every bin count still enters it.  A stronger compression is
possible: only $d$ disjoint aggregated count channels need to be retained.
The groups below are selected from the known candidate histograms; they do
not depend on the unknown mixture parameter.

\begin{theorem}[Minimal number of count channels]\label{mix:thm:counters}
For every sufficiently large $T$ and every allowed microscopic partition,
one can choose $d$ disjoint collections of bins
$\mathcal A_1,\ldots,\mathcal A_d$ such that observing only
\[
 Z_k=\sum_{B\in\mathcal A_k}Y_B,\qquad 1\le k\le d,
\tag{MF14}
\]
still gives squared-error risk
$O(\min\{1,(t\mathcal L_T)^{-1}\})$ uniformly on $[0,R]^d$.
These are independent Poisson counters.  The groups may unite bins at
widely separated heights, and all unselected bin counts are discarded.

Conversely, fewer than $d$ fixed disjoint aggregated Poisson counters
cannot identify an unrestricted $d$-dimensional interior box of mixture
parameters, even with arbitrarily long observation time.  This necessity
statement concerns ordinary count channels, not arbitrary real-valued
encodings of an entire dataset.
\end{theorem}

\begin{proof}
Write $p_B=v_B/s_B$, so $p_B$ lies in the probability simplex and
$G_T=\sum_Bs_Bp_Bp_B^{\mathsf T}$.  Let
$G_T\succeq c\mathcal L_T I_d$ and
$\sum_Bs_B\le C_0\mathcal L_T$, with constants from the frame theorem.
Partition the simplex into $J<\infty$ measurable cells of diameter at
most $\delta$, where $C_0\delta^2\le c/4$.  Both $J$ and $\delta$ are
fixed independently of $T$ and of the parameter.

For each nonempty cell $k$, aggregate its bins and put
\[
 S_k=\sum_{p_B\in k}s_B,\qquad
 \bar p_k=S_k^{-1}\sum_{p_B\in k}s_Bp_B,
 \qquad G_T^{\rm cell}=\sum_k S_k\bar p_k\bar p_k^{\mathsf T}.
\]
The exact within-cell covariance identity is
\[
 G_T-G_T^{\rm cell}
 =\sum_k\sum_{p_B\in k}s_B
          (p_B-\bar p_k)(p_B-\bar p_k)^{\mathsf T}
 \preceq C_0\delta^2\mathcal L_T I_d.
\tag{MF15}
\]
Thus $G_T^{\rm cell}\succeq(3c/4)\mathcal L_T I_d$.
Remove cells with $S_k<\eta\mathcal L_T$, where $\eta=c/(4J)$.
Since $\|\bar p_k\|_2\le1$, this removes a matrix of norm at most
$J\eta\mathcal L_T$, leaving a matrix $G_T^{\rm heavy}$ satisfying
\[
 G_T^{\rm heavy}\succeq(c/2)\mathcal L_T I_d,
 \qquad \eta\mathcal L_T\le S_k\le C_0\mathcal L_T
 \quad\hbox{for every remaining cell}.
\]

By Cauchy--Binet,
\[
 \det G_T^{\rm heavy}
 =\sum_{I:\,|I|=d}\left(\prod_{k\in I}S_k\right)
                 \det(P_I)^2,
\]
where $P_I$ has the rows $\bar p_k^{\mathsf T}$, $k\in I$.
Consequently some $d$ remaining cells satisfy
\[
 |\det P_I|\ge
 \kappa:=\left[\binom Jd^{-1}
                   \left(\frac{c}{2C_0}\right)^d\right]^{1/2}>0.
\tag{MF16}
\]
There are at least $d$ remaining cells because their matrix is positive
definite.  Every row has norm at most one, so
$\|P_I^{-1}\|_2\le d^{(d-1)/2}/\kappa$.

Retain precisely these $d$ cell counts.  They are independent with
\[
 Z_k\sim\operatorname{Pois}(tS_k\bar p_k^{\mathsf T}\theta).
\]
The estimator
\[
 \widehat\theta^{\rm count}
 =P_I^{-1}\left(\frac{Z_k}{tS_k}\right)_{k\in I}
\tag{MF17}
\]
is unbiased.  Its covariance is bounded above by
\[
 \frac{R}{t\eta\mathcal L_T}
                P_I^{-1}(P_I^{-1})^{\mathsf T}
 \preceq\frac{R d^{d-1}}{t\eta\kappa^2\mathcal L_T}I_d.
\]
Clipping to the box proves the claimed risk.  If
$A=\|P_I^{-1}\|_{\infty\to\infty}$, the usual Poisson Bernstein bound
also yields
\[
 \Pr\{\|\widehat\theta^{\rm count}-\theta\|_\infty\ge u\}
 \le2d\exp\left\{-\frac{t\eta\mathcal L_T(u/A)^2}
                              {2R+2u/(3A)}\right\}.
\]

Finally, $k<d$ disjoint count channels have a mean vector $tA\theta$ for
a fixed $k\times d$ matrix $A$.  Choose $b\ne0$ in its kernel and two
nearby interior points $\theta$ and $\theta+\varepsilon b$.  Their entire
count vectors have the same law at every $t$.  Thus such a channel system
is not identifiable on any open parameter box.
\end{proof}

This compression theorem uses only the frame bounds and nonnegative
intensity vectors.  The finite-dimensional Poisson argument is therefore
applicable to any family satisfying the same bounds.  The arithmetic
work is what proves those bounds despite shared and clustered zeros.

\begin{corollary}[Known total intensity]\label{mix:cor:normalized}
If $\boldsymbol1^{\mathsf T}\theta$ is known, then $d-1$ disjoint
aggregated counters suffice, with the same risk order on any compact
interior portion of that affine parameter simplex.  Fewer than $d-1$
cannot identify its relative interior.
\end{corollary}

\begin{proof}
Let $P$ be the well-conditioned matrix selected above and write
$\boldsymbol1^{\mathsf T}=c^{\mathsf T}P$.  Since every row of $P$ has
sum one, $\sum_jc_j=d$, so some $c_j\ge1$.  Replace that row of $P$ by
$\boldsymbol1^{\mathsf T}$.  The determinant of the resulting matrix has
absolute value $|c_j\det P|\ge\kappa$, and its row norms are at most
$\sqrt d$.  Use the known total in this row and retain only the other
$d-1$ counters.  The previous covariance argument applies with a zero
variance in the known-total coordinate.  Necessity is the kernel
argument restricted to the $(d-1)$-dimensional tangent hyperplane.
\end{proof}

\section{Why bounded microscopic resolution is the sharp regime}

\begin{theorem}[Collapse at growing bin width]\label{mix:thm:collapse}
Suppose $d\ge2$ and $K(T)\to\infty$.  There are partitions into intervals
of diameter at most $K(T)/\log T$ such that
\[
 \sup_{\substack{\|b\|_2=1\\\boldsymbol1^{\mathsf T}b=0}}
       b^{\mathsf T}G_Tb=o(\mathcal L_T),
 \qquad
 \boldsymbol1^{\mathsf T}G_T\boldsymbol1
       \asymp\mathcal L_T.
\tag{MF13}
\]
Consequently $G_T$ has one eigenvalue of order $\mathcal L_T$ and all
other eigenvalues $o(\mathcal L_T)$.  The same spectral conclusion holds
for $\mathcal I_T(\theta)/t$ uniformly on any fixed interior parameter box.
\end{theorem}

\begin{proof}
Put $m=\min\{K(T),\log T/\log\log T\}$ and $h=m/\log T$.  Partition the
band into intervals of width $h$, with a common origin selected uniformly
in $[0,h)$.  For every $j$, the expected sum over complete cells of
$|N_j(B)-N_1(B)|$ is
\[
 \frac1h\int_T^{2T-h}|C_j(u,h)-C_1(u,h)|du
 =O\left(\mathcal L_T\frac{\sqrt{\log m}}m\right)
\]
by (MF3).  The two boundary cells contribute $O(\log T)$ by the
unit-interval zero bound.  Choose a single origin for which the sum of
these discrepancies over all $j$ is at most its average.  Then
\[
 E_T:=\sum_B\sum_j|(v_B)_j-(v_B)_1|=o(\mathcal L_T).
\]
For $\boldsymbol1^{\mathsf T}b=0$, $\|b\|_2=1$,
$|b^{\mathsf T}v_B|\le\sum_j|(v_B)_j-(v_B)_1|$ and also
$|b^{\mathsf T}v_B|\le s_B$.  Hence
$\sum_B(b^{\mathsf T}v_B)^2/s_B\le E_T=o(\mathcal L_T)$ uniformly
in $b$.  Meanwhile
$\boldsymbol1^{\mathsf T}G_T\boldsymbol1=\sum_Bs_B\asymp\mathcal L_T$.
The eigenvalue assertion follows from the min--max principle and the
general upper bound.  Equation (MF11) transfers it to Fisher information.
\end{proof}

Thus a bounded number of mean zero spacings preserves all mixture
directions up to constant factors.  A growing number can preserve the
overall intensity while losing information at the leading $T\log T$ scale
about every balanced change of composition. Information of smaller order
can still diverge; no impossibility of consistent estimation follows from
this assertion alone. This is a statement about existence of
unfavorable growing-width partitions, not about every such partition.

\section{Application to the symmetric Li processes}

Choose primitive characters $\psi_1,\ldots,\psi_d$ with distinct
conjugacy orbits.  For $\gamma>0$ put
\[
 a_j(\gamma)=m_{\psi_j}(\gamma)+m_{\bar\psi_j}(\gamma).
\]
The symmetric Li process of the repository has rate $a_j(\gamma)/2$ at
each of the two jumps
$\pm2\arctan(1/(2\gamma))$.  For a real character this counts its
multiplicity twice in $a_j$, as required.  Summing the two signs gives
rate $a_j(\gamma)$; the sign is an independent fair mark and carries no
parameter information.  Group these magnitudes by ordinate bins and
use their $a$-count vectors in (MF6).

The frame theorem applies after the injective linear change from the
distinct primitive character family to its disjoint conjugacy-orbit sums.
For example, a coefficient vector on the orbit family expands to the
same coefficient on each member of a nonreal pair and to twice that
coefficient on a real member.  Its Euclidean norm is uniformly comparable
to the original norm because the family is fixed.  Every folded component
has the same leading zero count $(1/\pi)T\log T$, so the collapse theorem
uses the same balanced subspace.

Consequently the mixture estimator, confidence bound, minimax theorem
and sharp coarsening transition apply to intensity superpositions of the
symmetric Li processes.  Dirichlet GRH is used only to identify these
ordinate-built processes with the semigroups defined through the real
Li coefficients.  The histogram frame and its statistical consequences
for the ordinate-built point processes are unconditional.

The angular derivative is
$d[2\arctan(1/(2\gamma))]/d\gamma=-1/(\gamma^2+1/4)$.
On $(T,2T]$, bounded microscopic ordinate bins correspond to angle bins
of width $O(T^{-2}/\log T)$.  The observation model assumes calibrated
candidate zero histograms and exact bin membership.  It makes no claim
that noisy positions can be converted to those counts at the same cost,
and it does not assert an efficient uniform algorithm for constructing
the histograms or explicit constants uniform in conductor.

% Suggested bibliography entry:
% \bibitem{BP1998} E. Bombieri and A. Perelli, Distinct zeros of L-functions,
% Acta Arith. 83 (1998), 271--281, doi:10.4064/aa-83-3-271-281.
